% Unit 2 self-study -- "I do / You do" ladder for CED 2.1-2.7.
% Eleven ladders, forty-four questions.
%
% Disjoint from u02-frq, which uses the SCN- Lewis/formal-charge question,
% the NaCl/HCl/Cl2 bonding continuum, the Cl2-vs-Br2 bond energy comparison
% and ethanoic-acid hybridisation. This document uses K-F/C-O pairs, the
% LiF-NaF-KF lattice series, C-C/C=C/C#C bond order, and CO2/NO3-/SO4^2-
% for Lewis work instead.
\documentclass{apchem-worksheet}

\apcunit{2}
\apcunittitle{Compound Structure and Properties}
\apcdoctitle{Self-Study \textbullet\ CED 2.1--2.7, I~do / You~do}
\apcsubtitle{Eleven ladders \textbullet\ four YOUR TURN questions each \textbullet\ work alone, check after all four}

\begin{document}
\apctitleblock

\begin{apcnote}[How to use these notes --- read this first]
Each skill is a \textbf{ladder}: a worked example in the solid-framed box,
then \textbf{four} YOUR TURN questions in the dashed box --- same skill,
new substances, no help. Work all four before checking anything.

\textbf{The habit that decides this unit.} Almost every question here is
answered in two moves: \emph{count the electrons}, then \emph{reason about
the attraction}. Counting means valence electrons, domains, or charges.
Reasoning means \textbf{Coulomb's law} --- the product of the charges over
the distance between them.

An explanation that says a species ``wants a full octet'' or ``is more
stable'' earns \textbf{zero} on the AP exam. The octet rule tells you
\emph{what} the structure is; it never tells you \emph{why} a property
follows.

\textbf{One scope note.} For a central atom with more than four electron
domains the CED asks for the \textbf{shape only}. You will not be asked
for an \ce{sp^3d} or \ce{sp^3d^2} label, and none appears here.
\end{apcnote}

% =====================================================================
\section*{\sffamily Ladder 1 \textbullet\ Bond type from electronegativity}
\ced{2.1}

Bonding is a continuum, not three boxes. The electronegativity difference
$\Delta$EN tells you where on it a bond sits: near zero, the pair is
shared equally; large, the electron is effectively transferred.

\begin{workedexample}[I do: placing three bonds on the continuum]
Classify \ce{K-F}, \ce{C-O} and \ce{H-H} using
EN: H 2.1, C 2.5, O 3.5, F 4.0, K 0.8.

\textbf{\ce{K-F}:} $\Delta$EN $= 4.0 - 0.8 = 3.2$. Very large --- the
electron is transferred, giving \ce{K+} and \ce{F-}. \textbf{Ionic.}

\textbf{\ce{C-O}:} $\Delta$EN $= 3.5 - 2.5 = 1.0$. Moderate --- the pair
is shared but pulled toward oxygen, leaving $\delta+$ on C and $\delta-$
on O. \textbf{Polar covalent.}

\textbf{\ce{H-H}:} $\Delta$EN $= 0$. Identical atoms share equally.
\textbf{Nonpolar covalent.}

\textbf{Read the boundaries as guides, not laws.} Roughly: below
\num{0.4} nonpolar, \num{0.4}--\num{1.7} polar, above \num{1.7} ionic. A
bond at \num{1.6} is not categorically different from one at \num{1.8};
the character shifts gradually.
\end{workedexample}

\begin{yourturn}[YOUR TURN 1 --- four questions]
\begin{enumerate}[leftmargin=*,itemsep=0.5em]
  \item Classify \ce{Na-Br} (EN Na 0.9, Br 2.8): \workspace[1.4cm]
  \item Classify \ce{N-H} (EN N 3.0, H 2.1): \workspace[1.4cm]
  \item In \ce{P-Cl} (EN P 2.1, Cl 3.0), which atom carries $\delta-$, and
        why? \workspace[1.6cm]
  \item Two bonds have $\Delta$EN of \num{1.6} and \num{1.8}. Explain why
        calling one ``covalent'' and the other ``ionic'' is misleading.
        \workspace[1.8cm]
\end{enumerate}
\selfcheck{(a) ionic, $\Delta$EN \num{1.9} \quad (b) polar covalent,
\num{0.9} \quad (c) Cl, it is more electronegative \quad (d) the
continuum has no sharp boundary}
\keyonly{\begin{apcanswer}(a) $\Delta$EN $= 2.8 - 0.9 = \num{1.9}$, above
\num{1.7}: \textbf{ionic}.
(b) $\Delta$EN $= 3.0 - 2.1 = \num{0.9}$: \textbf{polar covalent}.
(c) \textbf{Chlorine}. It is the more electronegative atom, so it draws
the shared pair closer and gains the partial negative charge; phosphorus
is left $\delta+$.
(d) Because bonding is a \textbf{continuum}. A bond at \num{1.6} and one
at \num{1.8} have almost identical character --- both are strongly polar
with substantial charge separation. The \num{1.7} line is a convenient
convention, not a physical boundary where the nature of the bond suddenly
changes.\end{apcanswer}}
\end{yourturn}

% =====================================================================
\section*{\sffamily Ladder 2 \textbullet\ Coulomb's law and lattice energy}
\ced{2.2}

Lattice energy is the energy released when gaseous ions come together into
a crystal. It is Coulombic, so only two things matter: the
\textbf{product of the charges} and the \textbf{distance} between ion
centres.
\[ \text{attraction} \;\propto\; \frac{q_1 q_2}{d} \]

\begin{workedexample}[I do: one factor at a time]
Rank \ce{LiF}, \ce{NaF} and \ce{KF} by lattice energy, then compare
\ce{NaF} with \ce{MgO}.

\textbf{Within the fluoride series,} every compound pairs a $1+$ with a
$1-$, so the charge product is the same and cannot decide anything.
\emph{Distance} decides. Cation radii rise \ce{Li+} (76) $<$ \ce{Na+}
(102) $<$ \ce{K+} (138)~pm, so $d$ rises and the attraction falls:
\[ \ce{LiF} > \ce{NaF} > \ce{KF} \quad (1030 > 910 > 808~\text{kJ/mol}) \]

\textbf{\ce{NaF} against \ce{MgO}:} now the \emph{charges} change.
\ce{NaF} is $(+1)(-1) = 1$; \ce{MgO} is $(+2)(-2) = 4$ --- four times
the charge product. The ions are also smaller
(\SI{212}{\pico\metre} against \SI{235}{\pico\metre}), so $d$ is shorter
too. Both factors point the same way, and the gap is enormous:
\SI{3795}{} against \SI{910}{\kilo\joule\per\mole}.

\textbf{The method:} check charge first. If the charges match, distance
decides. If they differ, charge almost always dominates.
\end{workedexample}

\begin{yourturn}[YOUR TURN 2 --- four questions]
\begin{enumerate}[leftmargin=*,itemsep=0.5em]
  \item Larger lattice energy, \ce{NaCl} or \ce{NaF}? Why?
        \workspace[1.6cm]
  \item Larger lattice energy, \ce{CaO} or \ce{KF}? Why?
        \workspace[1.8cm]
  \item \ce{MgO} melts far higher than \ce{NaF}. Connect this to lattice
        energy. \workspace[1.8cm]
  \item Explain why comparing \ce{MgO} with \ce{NaCl} does \emph{not}
        isolate the effect of charge. \workspace[1.8cm]
\end{enumerate}
\selfcheck{(a) \ce{NaF} --- \ce{F-} smaller, shorter $d$ \quad
(b) \ce{CaO} --- charge product 4 vs 1 \quad (c) stronger attractions need
more energy to break \quad (d) charge \emph{and} size both differ}
\keyonly{\begin{apcanswer}(a) \textbf{\ce{NaF}}. Same charges, so distance
decides: \ce{F-} (133~pm) is smaller than \ce{Cl-} (181~pm), so $d$ is
shorter and the attraction stronger.
(b) \textbf{\ce{CaO}}. Charge product $(+2)(-2) = 4$ against \ce{KF}'s
$(+1)(-1) = 1$. Charge dominates, and \ce{CaO}'s ions are smaller too.
(c) Melting means giving ions enough energy to break out of the lattice.
A larger lattice energy means stronger electrostatic attractions holding
them, so more energy --- a higher temperature --- is needed.
(d) Because \textbf{both} variables change at once: \ce{MgO} has the
larger charge product \emph{and} the shorter internuclear distance. Any
difference cannot be attributed to charge alone. To isolate charge you
would compare species of similar size but different charge.\end{apcanswer}}
\end{yourturn}

% =====================================================================
\section*{\sffamily Ladder 3 \textbullet\ Bond order, length and energy}
\ced{2.2}

Within the same pair of atoms, more shared pairs means a shorter, stronger
bond. Bond order runs with energy and against length.

\begin{workedexample}[I do: a carbon--carbon series]
Rank \ce{C-C}, \ce{C=C} and \ce{C#C} by bond length and by bond energy.

Bond energies: 347, 614, 839 \si{\kilo\joule\per\mole}.

\textbf{Energy rises with bond order:} \ce{C-C} $<$ \ce{C=C} $<$
\ce{C#C}. Each extra shared pair adds attraction between the two nuclei
and the electron density between them.

\textbf{Length falls as order rises:} \ce{C#C} $<$ \ce{C=C} $<$ \ce{C-C}.
More shared electron density pulls the nuclei closer together.

\textbf{Note the energies are not simple multiples} --- 614 is not twice
347, and 839 is not three times it. The second and third pairs are pi
bonds, whose side-on overlap is less effective than the first bond's
end-on overlap.
\end{workedexample}

\begin{yourturn}[YOUR TURN 3 --- four questions]
\begin{enumerate}[leftmargin=*,itemsep=0.5em]
  \item Rank \ce{N-N} (163), \ce{N=N} (418), \ce{N#N} (941) by bond
        length, shortest first. \workspace[1.4cm]
  \item \ce{H-F} 567, \ce{H-Cl} 431, \ce{H-Br} 366, \ce{H-I} 299
        \si{\kilo\joule\per\mole}. Explain the trend.
        \workspace[2.2cm]
  \item Why is \ce{N2} so unreactive? \workspace[1.8cm]
  \item In \ce{O2} (495) and \ce{F2} (154), the bond order falls from 2 to
        1. Is that the whole explanation for the energy difference?
        \workspace[1.8cm]
\end{enumerate}
\selfcheck{(a) \ce{N#N} $<$ \ce{N=N} $<$ \ce{N-N} \quad (b) larger
halogen, longer bond, weaker attraction \quad (c) the triple bond is
extremely strong \quad (d) largely, but lone-pair repulsion in \ce{F2}
also weakens it}
\keyonly{\begin{apcanswer}(a) \ce{N#N} $<$ \ce{N=N} $<$ \ce{N-N}. Higher
bond order draws the nuclei closer.
(b) Down the halogen group the atom gets \textbf{larger}, so the
\ce{H-X} bond is longer and the shared pair sits further from both
nuclei. The Coulombic attraction holding the bond is correspondingly
weaker, so the bond energy falls steadily.
(c) Its \ce{N#N} triple bond is one of the strongest known
(\SI{941}{\kilo\joule\per\mole}). Any reaction must begin by breaking it,
which demands an enormous energy input, so nitrogen is kinetically
sluggish even where reaction is thermodynamically favourable.
(d) \textbf{Largely, but not entirely.} Bond order accounts for most of
it. \ce{F2} is also anomalously weak for a single bond because the two
small fluorine atoms are close together and their \textbf{lone pairs
repel} strongly, further destabilizing the bond --- which is why
\ce{F-F} (154) is weaker even than \ce{Cl-Cl} (243).\end{apcanswer}}
\end{yourturn}

% =====================================================================
\section*{\sffamily Ladder 4 \textbullet\ Properties of ionic solids}
\ced{2.3}

Every property of an ionic solid traces back to the same fact: charged
particles locked in a rigid lattice, held by strong electrostatic
attractions.

\begin{workedexample}[I do: four properties, one cause]
Explain why \ce{NaCl} is hard, brittle, high-melting, and conducts only
when molten or dissolved.

\textbf{Hard and high-melting:} each ion is attracted to all its
oppositely charged neighbours throughout the lattice. Deforming or melting
the solid means overcoming many strong Coulombic attractions at once.

\textbf{Brittle:} a blow displaces one plane of ions by a single lattice
position. Ions of \emph{like} charge are then brought face to face,
attraction along that plane becomes strong repulsion, and the crystal
splits cleanly.

\textbf{Conductivity:} conduction requires mobile charges. In the solid
the ions are fixed --- charges are present but cannot move. Melting or
dissolving frees them, and the liquid or solution conducts.

\textbf{The contrast that makes the point:} a metal, struck the same way,
bends. Its delocalized electrons flow with the displaced layer and keep
binding it, so no like-charge repulsion ever arises.
\end{workedexample}

\begin{yourturn}[YOUR TURN 4 --- four questions]
\begin{enumerate}[leftmargin=*,itemsep=0.5em]
  \item Why does solid \ce{KBr} not conduct, while molten \ce{KBr} does?
        \workspace[1.6cm]
  \item \ce{MgO} is harder than \ce{NaCl}. Explain Coulombically.
        \workspace[1.8cm]
  \item Why does an ionic crystal shatter rather than bend?
        \workspace[1.8cm]
  \item A solid dissolves in water and the solution conducts. Does that
        prove the solid was ionic? \workspace[1.8cm]
\end{enumerate}
\selfcheck{(a) ions fixed vs mobile \quad (b) charge product 4 vs 1 \quad
(c) like charges brought adjacent \quad (d) no --- a molecular acid also
ionizes}
\keyonly{\begin{apcanswer}(a) Conduction needs \emph{mobile} charge
carriers. In solid \ce{KBr} the ions sit in fixed lattice positions;
melting frees them to migrate, so the melt conducts.
(b) \ce{MgO} pairs $2+$ with $2-$, a charge product of 4, against
\ce{NaCl}'s 1; its ions are also smaller, so $d$ is shorter. Much stronger
attractions make the lattice far harder to deform.
(c) A blow shifts one layer by one ion position, bringing like charges
adjacent. Attraction along that plane is replaced by strong repulsion,
which drives the layers apart and splits the crystal.
(d) \textbf{No.} A molecular substance that \emph{ionizes} in water --- HCl,
for instance --- also gives a conducting solution despite being covalent
as a pure substance. The stronger evidence for an ionic solid is that it
conducts when \textbf{molten}, since that requires ions to have been
present in the solid itself.\end{apcanswer}}
\end{yourturn}

% =====================================================================
\section*{\sffamily Ladder 5 \textbullet\ Metals and alloys}
\ced{2.4}

A metal is a lattice of cations in a sea of delocalized electrons. Every
metallic property follows from those electrons being free to move.

\begin{workedexample}[I do: conductivity, malleability, alloys]
Explain metallic conduction and malleability, and why adding a second
element hardens the metal.

\textbf{Conduction:} the valence electrons belong to the whole lattice,
not to individual atoms. Under a potential difference they drift through
the solid, carrying charge.

\textbf{Malleability:} when layers of cations slide over one another, the
electron sea flows with them and continues to bind the lattice. Nothing
brings like charges together, so the metal deforms instead of fracturing.

\textbf{Alloying:} atoms of a different size disrupt the regular layers.
Layers can no longer slide smoothly past one another, so the alloy resists
deformation --- it is harder and stronger, though usually less ductile.
\end{workedexample}

\begin{yourturn}[YOUR TURN 5 --- four questions]
\begin{enumerate}[leftmargin=*,itemsep=0.5em]
  \item Why do metals conduct heat as well as electricity?
        \workspace[1.6cm]
  \item Brass is copper with zinc. Why is it harder than pure copper?
        \workspace[1.8cm]
  \item Why is an interstitial alloy such as steel usually harder than a
        substitutional one? \workspace[1.8cm]
  \item Both metals and molten ionic solids conduct. State how the
        \emph{carriers} differ. \workspace[1.8cm]
\end{enumerate}
\selfcheck{(a) mobile electrons carry kinetic energy \quad (b) zinc atoms
disrupt the layers \quad (c) small atoms wedge into gaps and pin the
layers \quad (d) electrons vs ions}
\keyonly{\begin{apcanswer}(a) The same delocalized electrons carry
\emph{kinetic} energy as well as charge. Heated at one end, they move
faster and transfer that energy rapidly through the lattice by collision.
(b) Zinc atoms differ in size from copper, disrupting the regular layers
so that they cannot slide over one another as easily.
(c) An interstitial alloy adds \emph{small} atoms (carbon in iron) that
wedge into the gaps between the metal atoms, pinning the layers directly.
A substitutional alloy merely swaps in atoms of similar size, which
disrupts the packing less.
(d) In a metal the mobile carriers are \textbf{electrons}, and the cations
stay put. In a molten ionic solid the carriers are the \textbf{ions}
themselves, which physically migrate through the liquid.\end{apcanswer}}
\end{yourturn}

% =====================================================================
\section*{\sffamily Ladder 6 \textbullet\ Drawing Lewis structures}
\ced{2.5}

Count the valence electrons first --- adjusting for charge --- then build a
skeleton, complete the outer atoms' octets, and convert lone pairs to
multiple bonds only if the central atom is short.

\begin{workedexample}[I do: a neutral molecule and an anion]
Draw \ce{CO2} and \ce{CO3^2-}.

\textbf{\ce{CO2}.} Valence electrons: $4 + 2(6) = \mathbf{16}$.
Skeleton \ce{O-C-O} uses 4. Completing octets on both oxygens uses the
other 12 --- but carbon then has only 4. Convert one lone pair on each
oxygen into a bond: \ce{O=C=O}, with two lone pairs on each oxygen. Every
atom now has an octet, and the total is still 16. \checkmark

\textbf{\ce{CO3^2-}.} Valence electrons:
$4 + 3(6) + \mathbf{2} = \mathbf{24}$ --- the $2-$ charge adds two.
Carbon central, three oxygens around it. One \ce{C=O} and two \ce{C-O},
with the whole structure in brackets carrying $2-$.

\textbf{The single most common error} is forgetting to adjust for the
charge. An anion has \emph{more} electrons than the neutral atoms
supply; a cation has fewer.
\end{workedexample}

\begin{yourturn}[YOUR TURN 6 --- four questions]
\begin{enumerate}[leftmargin=*,itemsep=0.5em]
  \item Valence-electron count for \ce{NH3}: \workspace[1.4cm]
  \item Valence-electron count for \ce{NH4+}: \workspace[1.4cm]
  \item Valence-electron count for \ce{SO4^2-}: \workspace[1.4cm]
  \item \ce{PCl3} has 26 valence electrons. How many lone pairs are on the
        phosphorus? \workspace[1.8cm]
\end{enumerate}
\selfcheck{(a) 8 \quad (b) 8 \quad (c) 32 \quad (d) one}
\keyonly{\begin{apcanswer}(a) $5 + 3(1) = \mathbf{8}$.
(b) $5 + 4(1) - \mathbf{1} = \mathbf{8}$ --- the $+$ charge \emph{removes}
one electron.
(c) $6 + 4(6) + \mathbf{2} = \mathbf{32}$.
(d) Three \ce{P-Cl} bonds use 6 electrons; completing the three chlorine
octets uses 18 more, totalling 24. The remaining $26 - 24 = 2$ electrons
sit on phosphorus as \textbf{one lone pair} --- which is what makes
\ce{PCl3} pyramidal rather than planar.\end{apcanswer}}
\end{yourturn}

% =====================================================================
\section*{\sffamily Ladder 7 \textbullet\ Formal charge}
\ced{2.6}

\[ \text{formal charge} = \text{valence electrons} -
   \text{lone-pair electrons} - \tfrac{1}{2}(\text{bonding electrons}) \]
Formal charges must sum to the overall charge on the species. The best
structure keeps them small, and puts any negative charge on the most
electronegative atom.

\begin{workedexample}[I do: nitrate]
Assign formal charges in \ce{NO3-}, drawn with one \ce{N=O} and two
\ce{N-O}.

\textbf{Nitrogen:} 5 valence, no lone pairs, 8 bonding electrons:
$5 - 0 - 4 = \mathbf{+1}$.

\textbf{Doubly bonded O:} 6 valence, 4 lone-pair electrons, 4 bonding:
$6 - 4 - 2 = \mathbf{0}$.

\textbf{Each singly bonded O:} 6 valence, 6 lone-pair electrons, 2
bonding: $6 - 6 - 1 = \mathbf{-1}$.

\textbf{Check the sum:} $(+1) + 0 + (-1) + (-1) = -1$, matching the ion's
charge. \checkmark If your formal charges do not sum to the overall
charge, the structure is drawn wrong --- go back before doing anything
else.
\end{workedexample}

\begin{yourturn}[YOUR TURN 7 --- four questions]
\begin{enumerate}[leftmargin=*,itemsep=0.5em]
  \item Formal charge on N in \ce{NH4+} (4 bonds, no lone pairs):
        \workspace[1.4cm]
  \item Formal charge on each O in \ce{CO2} (each doubly bonded, two lone
        pairs): \workspace[1.4cm]
  \item In \ce{SO4^2-} drawn with four single bonds, the formal charge on
        S is $+2$ and on each O is $-1$. Verify the sum.
        \workspace[1.6cm]
  \item Two structures for \ce{OCN-} put the $-1$ on O in one and on N in
        the other. Which is preferred? \workspace[1.8cm]
\end{enumerate}
\selfcheck{(a) $+1$ \quad (b) 0 \quad (c) $+2 + 4(-1) = -2$ \checkmark
\quad (d) the one with $-1$ on O, the more electronegative atom}
\keyonly{\begin{apcanswer}(a) $5 - 0 - \tfrac{1}{2}(8) = \mathbf{+1}$.
(b) $6 - 4 - \tfrac{1}{2}(4) = \mathbf{0}$ for each oxygen; carbon is
also 0, so the molecule sums to 0. \checkmark
(c) $(+2) + 4(-1) = \mathbf{-2}$, which matches the ion's charge.
\checkmark
(d) The structure with the negative formal charge on \textbf{oxygen}.
Oxygen is more electronegative than nitrogen and so accommodates negative
charge better; the arrangement placing $-1$ on the more electronegative
atom is the larger contributor.\end{apcanswer}}
\end{yourturn}

% =====================================================================
\section*{\sffamily Ladder 8 \textbullet\ Resonance}
\ced{2.6}

When a double bond can be placed in more than one equivalent position, no
single Lewis structure is correct. The real species is the
\textbf{average} --- the electrons are delocalized. It does not flip
between structures.

\begin{workedexample}[I do: nitrate's three structures]
Explain why \ce{NO3-} needs three resonance structures, and predict its
bond lengths.

\textbf{Why three:} the double bond can go to any one of the three
oxygens, and the three choices are equivalent.

\textbf{What is real:} not one structure, and not rapid switching between
them. The extra bonding pair is \textbf{delocalized} over all three
\ce{N-O} positions simultaneously.

\textbf{Bond lengths:} all three are \textbf{identical}, and each is
intermediate between a normal \ce{N-O} single bond and a normal \ce{N=O}
double bond. Every position has the same bond order, about
$1\tfrac{1}{3}$.

\textbf{The experimental point:} measurement finds three equal bond
lengths. A structure with one short and two long bonds would predict
otherwise, which is exactly why the single Lewis structure is known to be
inadequate.
\end{workedexample}

\begin{yourturn}[YOUR TURN 8 --- four questions]
\begin{enumerate}[leftmargin=*,itemsep=0.5em]
  \item How many equivalent resonance structures does \ce{CO3^2-} have?
        \workspace[1.4cm]
  \item Predict the relative \ce{C-O} bond lengths in \ce{CO3^2-}.
        \workspace[1.6cm]
  \item \ce{O3} has two resonance structures. What does this predict about
        its two \ce{O-O} bonds? \workspace[1.6cm]
  \item Why is ``the ion rapidly switches between structures'' wrong?
        \workspace[1.8cm]
\end{enumerate}
\selfcheck{(a) three \quad (b) all equal, between single and double \quad
(c) identical, bond order 1.5 \quad (d) the electrons are delocalized, not
alternating}
\keyonly{\begin{apcanswer}(a) \textbf{Three} --- the double bond can sit on
any of the three oxygens.
(b) All three \textbf{identical}, each intermediate between a \ce{C-O}
single and a \ce{C=O} double bond; bond order about $1\tfrac{1}{3}$.
(c) The two \ce{O-O} bonds are \textbf{identical}, each of bond order
\num{1.5} --- intermediate between single and double. Measurement confirms
this.
(d) Because nothing is moving between structures. The resonance
structures are limitations of the \emph{Lewis notation}, not states the
molecule occupies. The real electron distribution is a single, unchanging,
delocalized one that no individual Lewis structure can
draw.\end{apcanswer}}
\end{yourturn}

% =====================================================================
\section*{\sffamily Ladder 9 \textbullet\ VSEPR: shapes and angles}
\ced{2.7}

Count the electron \textbf{domains} on the central atom --- each bond
(single, double or triple) and each lone pair is one domain. Domains repel
and spread as far apart as possible. The \emph{shape} is then named from
the positions of the \textbf{atoms} only.

\begin{workedexample}[I do: same domain count, three different shapes]
Give the shape and approximate bond angle for \ce{CH4}, \ce{NH3} and
\ce{H2O}.

All three have \textbf{four} domains, so all three are based on a
tetrahedron. What differs is how many domains are lone pairs.

\textbf{\ce{CH4}} --- 4 bonds, 0 lone pairs: \textbf{tetrahedral},
\SI{109.5}{\degree}.

\textbf{\ce{NH3}} --- 3 bonds, 1 lone pair: \textbf{trigonal pyramidal},
about \SI{107}{\degree}.

\textbf{\ce{H2O}} --- 2 bonds, 2 lone pairs: \textbf{bent}, about
\SI{104.5}{\degree}.

\textbf{Why the angle shrinks:} a lone pair is held by only one nucleus,
so it spreads out closer to the central atom and takes up more angular
space than a bonding pair. Each lone pair squeezes the remaining bonds
together, so the angle drops with every one added.
\end{workedexample}

\begin{yourturn}[YOUR TURN 9 --- four questions]
\begin{enumerate}[leftmargin=*,itemsep=0.5em]
  \item Shape and angle of \ce{BF3}: \workspace[1.4cm]
  \item Shape and angle of \ce{SO2}: \workspace[1.6cm]
  \item Shape of \ce{SF6} (six domains, no lone pairs):
        \workspace[1.4cm]
  \item \ce{CO2} has two double bonds yet only two domains. Explain.
        \workspace[1.8cm]
\end{enumerate}
\selfcheck{(a) trigonal planar, \SI{120}{\degree} \quad (b) bent, just
under \SI{120}{\degree} \quad (c) octahedral \quad (d) a multiple bond
counts as one domain}
\keyonly{\begin{apcanswer}(a) Three domains, no lone pair:
\textbf{trigonal planar}, \SI{120}{\degree}.
(b) Three domains, one lone pair: \textbf{bent}, a little under
\SI{120}{\degree} because the lone pair takes more space.
(c) \textbf{Octahedral}, with \SI{90}{\degree} between adjacent
fluorines. (The CED asks only for the shape here --- no orbital label is
required.)
(d) A domain is a \emph{region} of electron density, not a count of
pairs. All the electrons of a double bond lie between the same two nuclei,
so they occupy one region and repel other domains as a single unit. Carbon
therefore has two domains and \ce{CO2} is linear.\end{apcanswer}}
\end{yourturn}

% =====================================================================
\section*{\sffamily Ladder 10 \textbullet\ Hybridization and sigma/pi counting}
\ced{2.7}

The central atom mixes exactly as many atomic orbitals as it has domains:
2 domains \ce{sp}, 3 domains \ce{sp^2}, 4 domains \ce{sp^3}. Every single
bond is one sigma; every double bond is one sigma plus one pi; every
triple bond is one sigma plus two pi.

\begin{workedexample}[I do: hybridization and bond count together]
For \ce{HCN}, give the hybridization of carbon and count the sigma and pi
bonds.

\textbf{Domains on carbon:} one bond to H, one triple bond to N. The
triple bond is a single domain, so carbon has \textbf{two} domains ---
hence \ce{sp} hybridized, and the molecule is linear.

\textbf{Bond count:} \ce{H-C} is one sigma. \ce{C#N} is one sigma plus
\emph{two} pi. Total: \textbf{2 sigma, 2 pi}.

\textbf{The rule that never fails:} the \emph{first} bond between any two
atoms is sigma; every additional bond between that same pair is pi.
\end{workedexample}

\begin{yourturn}[YOUR TURN 10 --- four questions]
\begin{enumerate}[leftmargin=*,itemsep=0.5em]
  \item Hybridization of B in \ce{BF3}: \workspace[1.4cm]
  \item Hybridization of O in \ce{H2O}: \workspace[1.4cm]
  \item Sigma and pi bonds in \ce{C2H4} (ethene): \workspace[1.6cm]
  \item Why does the \ce{C=C} bond in ethene resist rotation while the
        \ce{C-C} bond in ethane does not? \workspace[2.2cm]
\end{enumerate}
\selfcheck{(a) \ce{sp^2} \quad (b) \ce{sp^3} \quad (c) five sigma, one pi
\quad (d) rotation would destroy the pi overlap}
\keyonly{\begin{apcanswer}(a) Three domains, no lone pair:
\textbf{\ce{sp^2}}.
(b) Four domains --- two bonds and \emph{two lone pairs}:
\textbf{\ce{sp^3}}. Counting only the bonds would wrongly give
\ce{sp^2}.
(c) Four \ce{C-H} sigma plus one \ce{C-C} sigma $=$ \textbf{five sigma};
the second component of the \ce{C=C} is \textbf{one pi}.
(d) Ethane's \ce{C-C} is a single sigma bond, and sigma overlap is
cylindrically symmetric about the bond axis --- twisting does not reduce
it, so rotation is nearly free. Ethene's double bond includes a pi bond
formed by two \emph{parallel} $p$ orbitals; rotating one end turns them
out of alignment, and at \SI{90}{\degree} the overlap is zero. Rotation
would mean breaking the pi bond, which costs far too much
energy.\end{apcanswer}}
\end{yourturn}

% =====================================================================
\section*{\sffamily Ladder 11 \textbullet\ Molecular polarity}
\ced{2.1} \ced{2.7}

A molecule is polar only if it has polar bonds \emph{and} a shape that
prevents their dipoles cancelling. Shape decides --- which is why this
ladder comes after VSEPR.

\begin{workedexample}[I do: same bonds, opposite verdicts]
\ce{CO2} and \ce{H2O} both contain polar bonds. Explain why only water is
polar.

\textbf{\ce{CO2} is linear.} The two \ce{C=O} bond dipoles are equal and
point in exactly opposite directions, so they cancel completely. Net
dipole zero: \textbf{nonpolar}, despite strongly polar bonds.

\textbf{\ce{H2O} is bent}, because oxygen's two lone pairs force it so.
The two \ce{O-H} dipoles no longer oppose each other and add to a net
dipole pointing toward the oxygen: \textbf{polar}.

\textbf{The method:} identify the polar bonds, then ask whether the shape
is symmetric enough for them to cancel. Never answer from the bonds alone.
\end{workedexample}

\begin{yourturn}[YOUR TURN 11 --- four questions]
\begin{enumerate}[leftmargin=*,itemsep=0.5em]
  \item Is \ce{CCl4} polar? Justify. \workspace[1.6cm]
  \item Is \ce{CH3Cl} polar? Justify. \workspace[1.6cm]
  \item Is \ce{NH3} polar? Justify. \workspace[1.6cm]
  \item \ce{BF3} and \ce{NH3} both have three bonds to the central atom,
        but only one is polar. Explain. \workspace[2.0cm]
\end{enumerate}
\selfcheck{(a) no --- tetrahedral, dipoles cancel \quad (b) yes ---
asymmetric \quad (c) yes --- pyramidal \quad (d) \ce{BF3} is planar and
symmetric; \ce{NH3} has a lone pair and is pyramidal}
\keyonly{\begin{apcanswer}(a) \textbf{Nonpolar.} The four \ce{C-Cl} bonds
are polar, but the tetrahedral arrangement is perfectly symmetric and the
four dipoles cancel.
(b) \textbf{Polar.} Replacing one chlorine with hydrogen destroys that
symmetry: the \ce{C-Cl} dipole is no longer balanced by an identical bond
opposite it, so a net dipole remains.
(c) \textbf{Polar.} Three \ce{N-H} dipoles plus a lone pair on a
\emph{pyramidal} nitrogen; the dipoles do not cancel and the net points
along the lone-pair axis.
(d) \ce{BF3} has three domains and \emph{no} lone pair, so it is
\textbf{trigonal planar} --- perfectly symmetric, and the three dipoles
cancel. \ce{NH3} has four domains, one being a lone pair, so it is
\textbf{pyramidal} --- unsymmetric, and the dipoles add. Same number of
bonds, different domain count, opposite result.\end{apcanswer}}
\end{yourturn}

% =====================================================================
\section*{\sffamily Where you stand}

Tick a ladder only if all four YOUR TURN questions were right first time.

\renewcommand{\arraystretch}{1.35}
\noindent\begin{tabular}{@{}c p{6.2cm} c >{\raggedright\arraybackslash}p{4.8cm}@{}}
\toprule
\textbf{First try?} & \textbf{Skill} & \textbf{Ladder} &
  \textbf{If not, re-read\ldots}\\
\midrule
$\square$ & bond type from $\Delta$EN & 1 & it is a continuum \\
$\square$ & lattice energy & 2 & charge first, then distance \\
$\square$ & bond order, length, energy & 3 & more pairs, shorter, stronger \\
$\square$ & ionic solid properties & 4 & fixed ions; like charges repel \\
$\square$ & metals and alloys & 5 & delocalized electron sea \\
$\square$ & Lewis structures & 6 & adjust the count for charge \\
$\square$ & formal charge & 7 & must sum to the ion charge \\
$\square$ & resonance & 8 & delocalized, not alternating \\
$\square$ & VSEPR shape and angle & 9 & name from atoms, not domains \\
$\square$ & hybridization, sigma/pi & 10 & one orbital per domain \\
$\square$ & molecular polarity & 11 & shape decides, not bonds \\
\bottomrule
\end{tabular}

\begin{apcnote}[Scoring yourself honestly]
11/11: move on to the unit worksheets and the free-response set.
8--10: solid --- redo the missed ladders tomorrow, from a blank page.
7 or fewer: the recurring root causes in this unit are exactly three ---
(1) forgetting to adjust the valence-electron count for an ion's charge,
(2) counting bonds instead of \emph{domains}, so lone pairs get ignored,
and (3) answering a ``why'' question with the octet rule instead of with
charge, distance and shielding. Fix those three and most of these ladders
fall together.
\end{apcnote}

\end{document}
